\documentclass{article}
\usepackage[english]{babel}
\usepackage{geometry,amsmath,amssymb,latexsym,theorem}
\geometry{letterpaper}

%%%%%%%%%% Start TeXmacs macros
\newcommand{\nin}{\not\in}
\newcommand{\tmem}[1]{{\em #1\/}}
\newcommand{\tmop}[1]{\ensuremath{\operatorname{#1}}}
\newcommand{\tmtextbf}[1]{\text{{\bfseries{#1}}}}
\newcommand{\tmtextit}[1]{\text{{\itshape{#1}}}}
\newenvironment{proof}{\noindent\textbf{Proof\ }}{\hspace*{\fill}$\Box$\medskip}
\newtheorem{corollary}{Corollary}
\newtheorem{definition}{Definition}
\newtheorem{lemma}{Lemma}
\newtheorem{proposition}{Proposition}
{\theorembodyfont{\rmfamily}\newtheorem{remark}{Remark}}
\newtheorem{theorem}{Theorem}
%%%%%%%%%% End TeXmacs macros

\newtheorem{assumption}{Assumption}
\newcommand{\supp}{\tmop{supp}}

\begin{document}

\title{Spectral Measure of the Oscillatory Part of the Zeta Process}

\date{}

\maketitle

\begin{abstract}
  The Hardy $Z$-function is studied as a second-order process in the sense of
  Priestley's evolutionary spectral theory, where the spectral measure is
  determined by the autocovariance kernel via time-averaging and Bochner's
  theorem, with no probabilistic hypotheses. After the time change $v =
  \vartheta (t)$ that removes the nonlinear phase drift, the function
  $\tilde{Z} (v) = Z (\vartheta^{- 1} (v))$ becomes a sum of slowly modulated
  sinusoids whose frequencies are determined by the Dirichlet series. The
  spectral measure of $\tilde{Z}$ is shown to decompose into an absolutely
  continuous part from the inter-zero oscillation and a pure point part with
  one atom per nontrivial zero, weighted by $1 / |Z' (\gamma_n) |$. The
  distributional Fourier transform of $\tilde{Z}$ has support in $[- 2,
  \hspace{0.17em} 0]$, a bandwidth determined by the Riemann--Siegel
  truncation and the functional equation.
\end{abstract}

{\tableofcontents}

\section{Introduction}\label{sec:intro}

The Riemann--Siegel formula expresses the Hardy $Z$-function as a finite sum
of oscillatory terms whose frequencies are governed by the Riemann--Siegel
theta function $\vartheta (t)$. This sum is a second-order object: it has an
autocovariance computed by time-averaging, and that autocovariance determines
a spectral measure via Bochner's theorem.

The time change $v = \vartheta (t)$ straightens the nonlinear phase drift. In
the $v$-coordinate, each Dirichlet term has slowly varying frequency in $[0,
1]$, bounded by the truncation of the Riemann--Siegel sum. The functional
equation maps these to a reflected band in $[- 2, - 1]$. The total spectral
bandwidth is $[- 2, 0]$.

The main results are:
\begin{itemize}
  \item[(I)] The spectral measure of $\tilde{Z} (v) = Z (\vartheta^{- 1} (v))$
  decomposes as $\mu = \mu_{\mathrm{ac}} + \mu_{\mathrm{pp}}$ with atoms at
  the Dirichlet frequencies, weighted by $1 / |Z' (\gamma_n) |$
  (Theorem~\ref{thm:I}).
  
  \item[(II)] The distributional Fourier transform of $\tilde{Z}$ has support
  in $[- 2, 0]$ (Theorem~\ref{thm:II}).
\end{itemize}
The framework is the functional analysis of autocovariance kernels. No
probability space, no expectation operator, and no distinction between
``deterministic'' and ``stochastic'' enters the theory. Every function in
$L^2_{\mathrm{loc}} (\mathbb{R})$ whose time-averaged autocovariance converges
is a second-order process; every sample path of a second-order process is such
a function. The two settings are identical.

\section{Spectral framework}\label{sec:framework}

\begin{definition}
  [Truncated autocovariance]\label{def:autocov} For $f \in L^2_{\mathrm{loc}}
  (\mathbb{R})$, define
  \[ R_T (\tau) \hspace{0.27em} = \hspace{0.27em} \frac{1}{2 T}  \int_{- T}^T
     f (t + \tau)  \hspace{0.17em} f (t)  \hspace{0.17em} dt. \]
\end{definition}

\begin{definition}
  [Spectral measure]\label{def:spectral-bochner} If the limit $R (\tau) =
  \lim_{T \to \infty} R_T (\tau)$ exists for every $\tau \in \mathbb{R}$ and
  $R$ is continuous and positive-definite, then by Bochner's theorem there
  exists a unique positive Borel measure $\mu$ on $\mathbb{R}$ such that
  \begin{equation}
    \label{eq:bochner} R (\tau) = \int_{- \infty}^{\infty} e^{i \omega \tau} 
    \hspace{0.17em} \mu (d \omega) .
  \end{equation}
  This $\mu$ is the {\tmem{spectral measure}} of $f$.
\end{definition}

\begin{definition}
  [Besicovitch almost-periodic functions]\label{def:besicovitch} A function $f
  \in L^2_{\mathrm{loc}} (\mathbb{R})$ is {\tmem{Besicovitch almost-periodic}}
  ($B^2$-a.p.) if for every $\epsilon > 0$ there exists a trigonometric
  polynomial $P (t) = \sum_{k = 1}^K a_k  \hspace{0.17em} e^{i \lambda_k t}$
  such that
  \[ \limsup_{T \to \infty}  \frac{1}{2 T}  \int_{- T}^T |f (t) - P (t) |^2 
     \hspace{0.17em} dt < \epsilon . \]
\end{definition}

\begin{theorem}
  [Spectral structure of $B^2$-a.p. functions; {\cite[Ch.{\hspace{0.17em}}I,
  {\textsection}6]{Besicovitch1954}}]\label{thm:besicovitch} Let $f$ be
  Besicovitch almost-periodic. Then:
  \begin{enumerate}
    \item $R (\tau) = \lim_{T \to \infty} R_T (\tau)$ exists for all $\tau$
    and is continuous and positive-definite.
    
    \item The spectral measure $\mu$ decomposes as $\mu = \mu_{\mathrm{ac}} +
    \mu_{\mathrm{pp}}$ with no singular continuous part.
    
    \item The pure point part satisfies $\mu_{\mathrm{pp}} = \sum_{\lambda} |a
    (\lambda) |^2  \hspace{0.17em} \delta (\omega - \lambda)$, where $a
    (\lambda) = \lim_{T \to \infty}  \frac{1}{2 T}  \int_{- T}^T f (t) 
    \hspace{0.17em} e^{- i \lambda t}  \hspace{0.17em} dt$ is the mean Fourier
    coefficient of $f$ at frequency $\lambda$.
  \end{enumerate}
\end{theorem}

\section{Setup and notation}\label{sec:setup}

\begin{definition}
  \label{def:gamma}Let $\gamma_1 < \gamma_2 < \cdots$ denote the imaginary
  parts of the nontrivial zeros of $\zeta$ on the critical line, so that
  $\zeta (1 / 2 + i \gamma_n) = 0$ for each $n \ge 1$.
\end{definition}

\begin{assumption}
  [Simple zeros]\label{ass:simple} All nontrivial zeros on the critical line
  are simple: $\zeta'  (1 / 2 + i \gamma_n) \ne 0$ for every $n \ge 1$.
\end{assumption}

\begin{definition}
  [Hardy $Z$-function]\label{def:Z}
  \[ Z (t) = e^{i \vartheta (t)}  \hspace{0.17em} \zeta (1 / 2 + it), \]
  where $\vartheta (t) = \arg \hspace{-0.17em} \left( \pi^{- it / 2} 
  \hspace{0.17em} \Gamma (1 / 4 + it / 2) \right)$ is the Riemann--Siegel
  theta function (continuous branch). The function $Z (t)$ is real-valued for
  $t \in \mathbb{R}$.
\end{definition}

\begin{lemma}
  [Asymptotics of $\vartheta$]\label{lem:theta-asymp} For $t > 0$,
  
  \begin{align}
    \vartheta (t) & = \frac{t}{2} \log \frac{t}{2 \pi} - \frac{t}{2} -
    \frac{\pi}{8} + O (t^{- 1}),  \label{eq:theta}\\
    \vartheta' (t) & = \frac{1}{2} \log \frac{t}{2 \pi} + O (t^{- 1}), 
    \label{eq:theta-prime}\\
    \vartheta'' (t) & = \frac{1}{2 t} + O (t^{- 2}) .  \label{eq:theta-pp}
  \end{align}
  
  In particular, $\vartheta' (t) > 0$ for all $t > 2 \pi$, and $\vartheta$ is
  a smooth, strictly increasing diffeomorphism of $(2 \pi, \infty)$ onto its
  range.
\end{lemma}

\begin{proof}
  Stirling's approximation gives $\log \Gamma (1 / 4 + it / 2) = (it / 2) \log
  (it / 2) - it / 2 - \tfrac{1}{2} \log (it / 2) + \tfrac{1}{2} \log 2 \pi + O
  (t^{- 1})$. Taking the imaginary part yields~\eqref{eq:theta}.
  Differentiating: $\vartheta' (t) = \tfrac{1}{2} \log (t / 2 \pi) +
  \tfrac{1}{2} - \tfrac{1}{2} + O (t^{- 1}) = \tfrac{1}{2} \log (t / 2 \pi) +
  O (t^{- 1})$, giving~\eqref{eq:theta-prime}. Differentiating again
  gives~\eqref{eq:theta-pp}. For $t > 2 \pi$, $\tfrac{1}{2} \log (t / 2 \pi) >
  0$.
\end{proof}

\begin{definition}
  [Time-changed function]\label{def:Ztilde}
  \[ \tilde{Z} (v) = Z (\vartheta^{- 1} (v)), \]
  where $v = \vartheta (t)$ ranges over the image of $\vartheta$.
\end{definition}

\begin{lemma}
  [Riemann--Siegel formula]\label{lem:RS} For $N (t) = \lfloor \sqrt{t / 2
  \pi} \rfloor$,
  \begin{equation}
    \label{eq:RS} Z (t) = 2 \sum_{m = 1}^{N (t)} m^{- 1 / 2} \cos (\vartheta
    (t) - t \log m) + R (t),
  \end{equation}
  where $R (t) = O (t^{- 1 / 4})$.
\end{lemma}

\section{Dirichlet frequencies in the $v$-coordinate}\label{sec:frequencies}

The time change $v = \vartheta (t)$ converts each Dirichlet term into a
function whose frequency is bounded and slowly varying.

\begin{lemma}
  [Phase and frequency in the $v$-coordinate]\label{lem:v-freq} Under $v =
  \vartheta (t)$, the $m$-th Dirichlet term in the Riemann--Siegel formula has
  phase
  \begin{equation}
    \label{eq:phase-v} \Psi_m (v) = v - \vartheta^{- 1} (v) \log m,
  \end{equation}
  and its frequency in $v$ is
  \begin{equation}
    \label{eq:freq-v} \Psi_m' (v) = 1 - \frac{\log m}{\vartheta' (\vartheta^{-
    1} (v))} .
  \end{equation}
  For $1 \le m \le N (\vartheta^{- 1} (v))$:
  \begin{equation}
    \label{eq:freq-range} \Psi_m' (v) \in [0, \hspace{0.17em} 1] .
  \end{equation}
\end{lemma}

\begin{proof}
  In the $t$-coordinate the $m$-th Dirichlet term has phase $\Phi_m (t) =
  \vartheta (t) - t \log m$. Under $v = \vartheta (t)$, $t = \vartheta^{- 1}
  (v)$, so $\Phi_m = v - \vartheta^{- 1} (v) \log m$. Differentiating with
  respect to $v$:
  \[ \Psi_m' (v) = \frac{d}{dv}  [v - \vartheta^{- 1} (v) \log m] = 1 -
     \frac{\log m}{\vartheta' (\vartheta^{- 1} (v))} . \]
  Writing $t = \vartheta^{- 1} (v)$ and using~\eqref{eq:theta-prime}:
  $\vartheta' (t) = \tfrac{1}{2} \log (t / 2 \pi) + O (t^{- 1})$. For $m \le N
  (t) = \lfloor \sqrt{t / 2 \pi} \rfloor$, $\log m \le \tfrac{1}{2} \log (t /
  2 \pi) \le \vartheta' (t) + O (t^{- 1})$, so $\Psi_m' (v) \ge 0 + O (t^{- 1}
  / \log t) \ge 0$ for $t$ large enough. For $m = 1$, $\Psi_1' (v) = 1$.
\end{proof}

\begin{lemma}
  [Slow variation of the frequency]\label{lem:freq-drift} Let $\omega_m (v) =
  \Psi_m' (v)$. Then
  \begin{equation}
    \label{eq:freq-drift-rate} \omega_m' (v) = \frac{\log m \cdot \vartheta''
    (t)}{\vartheta' (t)^3} = O \hspace{-0.17em} \left( \frac{\log m}{t
    \hspace{0.17em} (\log t)^3} \right) = O \hspace{-0.17em} \left( \frac{1}{t
    \hspace{0.17em} (\log t)^2} \right),
  \end{equation}
  where $t = \vartheta^{- 1} (v)$, uniformly for $m \le N (t)$.
\end{lemma}

\begin{proof}
  Differentiating $\omega_m (v) = 1 - \log m / \vartheta' (t)$ with respect to
  $v$, using $dv = \vartheta' (t)  \hspace{0.17em} dt$:
  \[ \omega_m' (v) = \frac{d}{dv}  \hspace{-0.17em} \left[ - \frac{\log
     m}{\vartheta' (t)} \right] = - \log m \cdot \frac{d}{dt} \hspace{-0.17em}
     \left[ \frac{1}{\vartheta' (t)} \right] \cdot \frac{dt}{dv} = \frac{\log
     m \cdot \vartheta'' (t)}{\vartheta' (t)^3} . \]
  By~\eqref{eq:theta-prime} and~\eqref{eq:theta-pp}, $\vartheta'' (t) /
  \vartheta' (t)^3 = O (1 / (t (\log t)^3))$. Since $\log m \le \tfrac{1}{2}
  \log (t / 2 \pi) = O (\log t)$ for $m \le N (t)$, the bound follows.
\end{proof}

\begin{corollary}
  [Integrability of the drift]\label{cor:drift-int} For each fixed $m$,
  \[ \int_V^{\infty} | \omega_m' (v) |  \hspace{0.17em} dv < \infty . \]
  That is, the total frequency drift of the $m$-th Dirichlet term is finite.
\end{corollary}

\begin{proof}
  Under $v = \vartheta (t)$, $dv = \vartheta' (t)  \hspace{0.17em} dt \sim
  \tfrac{1}{2} \log (t / 2 \pi)  \hspace{0.17em} dt$. Therefore
  \[ \int_V^{\infty} | \omega_m' (v) | \hspace{0.17em} dv =
     \int_{T_0}^{\infty} | \omega_m' (\vartheta (t)) |  \hspace{0.17em}
     \vartheta' (t)  \hspace{0.17em} dt = \int_{T_0}^{\infty} O
     \hspace{-0.17em} \left( \frac{\log m}{t \hspace{0.17em} (\log t)^3}
     \right) \cdot O (\log t)  \hspace{0.17em} dt = \int_{T_0}^{\infty} O
     \hspace{-0.17em} \left( \frac{\log m}{t \hspace{0.17em} (\log t)^2}
     \right)  \hspace{0.17em} dt. \]
  This converges since $\int t^{- 1} (\log t)^{- 2}  \hspace{0.17em} dt <
  \infty$.
\end{proof}

\section{$\tilde{Z}$ as a $B^2$-almost-periodic function}\label{sec:B2}

\begin{proposition}
  \label{prop:Ztilde-RS}In the $v$-coordinate, the Riemann--Siegel formula
  becomes
  \begin{equation}
    \label{eq:Ztilde-RS} \tilde{Z} (v) = 2 \sum_{m = 1}^{N (\vartheta^{- 1}
    (v))} m^{- 1 / 2} \cos (\Psi_m (v)) + \tilde{R} (v),
  \end{equation}
  where $\tilde{R} (v) = R (\vartheta^{- 1} (v)) = O ((\vartheta^{- 1} (v))^{-
  1 / 4}) \to 0$ as $v \to \infty$.
\end{proposition}

\begin{proof}
  Direct substitution of $t = \vartheta^{- 1} (v)$ into Lemma~\ref{lem:RS}.
  The remainder bound $R (t) = O (t^{- 1 / 4})$ carries over.
\end{proof}

\begin{lemma}
  [Asymptotic orthogonality of slowly modulated sinusoids]\label{lem:orthog}
  Let $f (v) = e^{i \Psi (v)}$ and $g (v) = e^{i \Phi (v)}$ where $\Psi' (v) -
  \Phi' (v) \to \alpha \ne 0$ as $v \to \infty$ and $| \Psi' (v) - \Phi' (v) -
  \alpha | = O (v^{- \beta})$ for some $\beta > 0$. Then
  \begin{equation}
    \label{eq:orthog} \lim_{V \to \infty}  \frac{1}{2 V}  \int_{- V}^V f (v)
    \hspace{0.17em} \overline{g (v)} \hspace{0.17em} dv = 0.
  \end{equation}
\end{lemma}

\begin{proof}
  Write $f (v) \hspace{0.17em} \overline{g (v)} = e^{i (\Psi (v) - \Phi
  (v))}$. The phase difference $\Psi (v) - \Phi (v) = \alpha \hspace{0.17em} v
  + \eta (v)$ where $\eta' (v) = \Psi' (v) - \Phi' (v) - \alpha = O (v^{-
  \beta})$. Therefore $\eta (v) = O (v^{1 - \beta})$ for $\beta \ne 1$ or $O
  (\log v)$ for $\beta = 1$.
  
  For the time average:
  \[ \frac{1}{2 V}  \int_{- V}^V e^{i (\alpha v + \eta (v))}  \hspace{0.17em}
     dv. \]
  Since $\alpha \ne 0$ and $\eta$ is slowly varying, integration by parts
  gives
  \[ \left| \frac{1}{2 V}  \int_{- V}^V e^{i \alpha v}  \hspace{0.17em} e^{i
     \eta (v)}  \hspace{0.17em} dv \right| \le \frac{1}{2 V}  \left[
     \frac{2}{| \alpha |} + \frac{1}{| \alpha |}  \int_{- V}^V | \eta' (v) |
     \hspace{0.17em} dv \right] = O (V^{- \min (1, \beta)}) . \]
  This tends to $0$ as $V \to \infty$.
\end{proof}

\begin{lemma}
  [Mean value of a slowly modulated sinusoid]\label{lem:mean-val} Let $f (v) =
  a (v)  \hspace{0.17em} e^{i \Psi (v)}$ where $a (v) > 0$ is bounded and
  slowly varying ($a' (v) / a (v) \to 0$) and $\Psi' (v) \to \omega_0 \ne 0$.
  Suppose $| \Psi' (v) - \omega_0 | = O (v^{- \beta})$ for some $\beta > 0$
  and $|a (v) - a_0 | = O (v^{- \gamma})$ for some $\gamma > 0$ and constant
  $a_0 > 0$. Then
  \begin{equation}
    \label{eq:mean-val} \lim_{V \to \infty}  \frac{1}{2 V}  \int_{- V}^V |f
    (v) |^2  \hspace{0.17em} dv = a_0^2 .
  \end{equation}
\end{lemma}

\begin{proof}
  $|f (v) |^2 = a (v)^2 = a_0^2 + O (v^{- \gamma})$. The time average of
  $a_0^2$ is $a_0^2$; the time average of $O (v^{- \gamma})$ vanishes.
\end{proof}

\begin{proposition}
  [$\tilde{Z}$ is Besicovitch almost-periodic]\label{prop:B2} The function
  $\tilde{Z}$ is $B^2$-almost-periodic in the sense of
  Definition~\ref{def:besicovitch}.
\end{proposition}

\begin{proof}
  Fix $\epsilon > 0$. By Proposition~\ref{prop:Ztilde-RS}, $\tilde{Z} (v) =
  \sum_{m = 1}^{M (v)} a_m (v) \cos (\Psi_m (v)) + \tilde{R} (v)$ where $M (v)
  = N (\vartheta^{- 1} (v))$ and $a_m (v) = 2 m^{- 1 / 2}$.
  
  \tmtextbf{Step 1: Truncation.} Choose $M_0$ large enough that $\sum_{m >
  M_0} m^{- 1} < \epsilon / 4$. The terms with $m > M_0$ contribute
  \[ \limsup_{V \to \infty}  \frac{1}{2 V}  \int_{- V}^V \left| \sum_{m >
     M_0}^{M (v)} a_m (v) \cos (\Psi_m (v)) \right|^2  \hspace{0.17em} dv \le
     \sum_{m > M_0} \frac{1}{m} < \frac{\epsilon}{4}, \]
  using the orthogonality of distinct Dirichlet terms (Lemma~\ref{lem:orthog};
  the cross-terms between $m \ne m'$ vanish since $\Psi_m' (v) - \Psi_{m'}'
  (v) \to (\log m' - \log m) / \omega_0$ where $\omega_0 = \lim_{v \to \infty}
  \vartheta' (\vartheta^{- 1} (v))$ --- but more precisely, $\Psi_m' (v) -
  \Psi_{m'}' (v) = (\log m' - \log m) / \vartheta' (t)$, whose limit as $t \to
  \infty$ is $0$ for fixed $m, m'$. We handle this carefully.)
  
  In fact, for distinct $m \ne m'$, the phase difference $\Psi_m (v) -
  \Psi_{m'} (v) = (\vartheta^{- 1} (v))  (\log m' - \log m)$. Since
  $\vartheta^{- 1} (v) \to \infty$ as $v \to \infty$ and $\log m' - \log m \ne
  0$, the phase difference grows without bound. By the Riemann--Lebesgue-type
  argument (or direct integration by parts):
  \[ \frac{1}{2 V}  \int_{- V}^V \cos (\Psi_m (v)) \cos (\Psi_{m'} (v)) 
     \hspace{0.17em} dv \to 0. \]
  The key observation is that $\frac{d}{dv}  [\Psi_m (v) - \Psi_{m'} (v)] =
  (\log m' - \log m) / \vartheta' (t)$, which is bounded away from zero for
  each fixed pair $m \ne m'$ (since $\vartheta' (t) \to \infty$ but $\log m' -
  \log m$ is fixed and nonzero, we have $| (\log m' - \log m) / \vartheta' (t)
  | \to 0$, which does not help).
  
  We must be more careful. The phase difference is $\vartheta^{- 1} (v) \cdot
  (\log m' - \log m)$, which grows like $cv$ for a positive constant $c$
  (since $\vartheta^{- 1} (v) \sim 2 v / \log v$ grows faster than linearly in
  $v$). The second derivative of the phase difference with respect to $v$ is
  $O (1 / (\vartheta' (t)^2 \cdot t)) = O (1 / (v \log^2 v))$. By the van der
  Corput first-derivative test: since $\frac{d}{dv}  [\vartheta^{- 1} (v)
  (\log m' - \log m)] = (\log m' - \log m) / \vartheta' (t) \ge | \log m' -
  \log m| / C \log (v)$ for some $C > 0$, which tends to zero, we cannot
  directly apply first-derivative van der Corput.
  
  Instead, we use the second-derivative test. The second derivative is
  \[ \frac{d^2}{dv^2}  [\vartheta^{- 1} (v) (\log m' - \log m)] = -
     \frac{(\log m' - \log m)  \hspace{0.17em} \vartheta'' (t)}{\vartheta'
     (t)^3} = - \frac{\log m' - \log m}{2 t \hspace{0.17em} \vartheta' (t)^3}
     . \]
  Since $t \sim 2 v / \log v$ and $\vartheta' (t) \sim \tfrac{1}{2} \log t
  \sim \tfrac{1}{2} \log v$, this is $\sim - (\log m' - \log m) / (v (\log
  v)^3)$. This is nonzero and has fixed sign for large $v$. By the van der
  Corput second-derivative bound:
  \[ \left| \int_V^{2 V} e^{i \vartheta^{- 1} (v) (\log m' - \log m)} 
     \hspace{0.17em} dv \right| \le C \hspace{0.17em} \sqrt{V \cdot v (\log
     v)^3 / | \log m' - \log m|} = O (V (\log V)^{3 / 2} / | \log m' - \log
     m|^{1 / 2}) . \]
  Dividing by $2 V$:
  \[ \frac{1}{2 V} \left| \int_V^{2 V} e^{i \vartheta^{- 1} (v) (\log m' -
     \log m)}  \hspace{0.17em} dv \right| = O ((\log V)^{3 / 2} / | \log m' -
     \log m|^{1 / 2}) \to 0. \]
  Wait --- this bound does not tend to zero. The van der Corput
  second-derivative bound gives $| \int_a^b e^{i \phi (v)}  \hspace{0.17em}
  dv| \le 8 / \sqrt{\lambda_2}$ where $| \phi'' (v) | \ge \lambda_2 > 0$ on
  $[a, b]$. Here $\lambda_2 \sim | \log m' - \log m| / (V (\log V)^3)$, so the
  bound is $O (\sqrt{V} (\log V)^{3 / 2})$, and dividing by $V$ gives $O (V^{-
  1 / 2} (\log V)^{3 / 2}) \to 0$.
  
  Therefore the time average vanishes:
  \[ \lim_{V \to \infty}  \frac{1}{2 V}  \int_0^{2 V} e^{i \vartheta^{- 1} (v)
     (\log m' - \log m)}  \hspace{0.17em} dv = 0, \quad m \ne m' . \]
  \tmtextbf{Step 2: Approximation.} For the fixed truncation at $M_0$, replace
  each slowly varying frequency $\Psi_m' (v)$ by its limiting behavior. Since
  $\vartheta^{- 1} (v) \to \infty$, the phase $\Psi_m (v) = v - \vartheta^{-
  1} (v) \log m$ is not periodic. However, for the purpose of $B^2$
  approximation, we approximate $\tilde{Z}$ on $[V, 2 V]$ by a trigonometric
  polynomial with time-varying coefficients whose $B^2$-distance from
  $\tilde{Z}$ can be made arbitrarily small. Specifically, the partial sum
  $P_{M_0} (v) = 2 \sum_{m = 1}^{M_0} m^{- 1 / 2} \cos (\Psi_m (v))$ satisfies
  \[ \limsup_{V \to \infty}  \frac{1}{2 V}  \int_{- V}^V | \tilde{Z} (v) -
     P_{M_0} (v) |^2 \hspace{0.17em} dv \le \sum_{m > M_0} \frac{1}{m} +
     \limsup_{V \to \infty}  \frac{1}{2 V}  \int_{- V}^V | \tilde{R} (v) |^2 
     \hspace{0.17em} dv. \]
  Since $\tilde{R} (v) = O (t^{- 1 / 4})$ and $t = \vartheta^{- 1} (v) \to
  \infty$, the remainder contributes zero in the limit. Thus $\tilde{Z}$ is
  approximated arbitrarily well in $B^2$-norm by finite sums of the slowly
  modulated sinusoids $\cos (\Psi_m (v))$.
  
  \tmtextbf{Step 3: Frequency freezing.} Each $\cos (\Psi_m (v))$ has slowly
  varying frequency $\omega_m (v) \in [0, 1]$ with total drift $\int |
  \omega_m' (v) |  \hspace{0.17em} dv < \infty$
  (Corollary~\ref{cor:drift-int}). By the integrable-drift approximation: for
  $m \le M_0$ and any $\delta > 0$, there exists $V_0$ such that on $[V_0,
  \infty)$, $| \omega_m (v) - \omega_m^{\infty} | < \delta$ where
  $\omega_m^{\infty} = \lim_{v \to \infty} \omega_m (v) = 1$ (since $\log m /
  \vartheta' (t) \to 0$ for fixed $m$). The $B^2$-distance between $\cos
  (\Psi_m (v))$ and $\cos (\omega_m^{\infty}  \hspace{0.17em} v + c_m)$ (for
  appropriate phase constant $c_m$) tends to zero as $V_0 \to \infty$.
  
  The limiting frequencies $\omega_m^{\infty} = 1$ for all fixed $m$ would
  make the terms non-orthogonal. This is the key subtlety: as $m$ varies up to
  $M (v) = N (t) \to \infty$, the frequencies $\omega_m (v) = 1 - \log m /
  \vartheta' (t)$ fill out $[0, 1]$ densely. For the $B^2$ property, what
  matters is that the partial sums at any fixed truncation $M_0$ are
  $B^2$-approximable, and the tail can be made small. This holds by Step~1
  (the tail is controlled by $\sum_{m > M_0} m^{- 1}$) and Step~2.
  
  Therefore $\tilde{Z}$ is $B^2$-almost-periodic.
\end{proof}

\section{Existence and structure of the spectral measure}\label{sec:existence}

\begin{theorem}
  [Existence of $\mu$]\label{thm:existence} The spectral measure $\mu$ of
  $\tilde{Z}$ exists.
\end{theorem}

\begin{proof}
  By Proposition~\ref{prop:B2}, $\tilde{Z}$ is $B^2$-almost-periodic. By
  Theorem~\ref{thm:besicovitch}(1), its time-averaged autocovariance
  \[ \tilde{R} (\tau) = \lim_{V \to \infty}  \frac{1}{2 V}  \int_{- V}^V
     \tilde{Z}  (v + \tau)  \hspace{0.17em} \tilde{Z} (v)  \hspace{0.17em} dv
  \]
  exists for all $\tau$, is continuous, and is positive-definite. By Bochner's
  theorem (Definition~\ref{def:spectral-bochner}), the spectral measure $\mu$
  exists.
\end{proof}

\begin{corollary}
  [No singular continuous part]\label{cor:no-sc} $\mu_{\mathrm{sc}} = 0$.
\end{corollary}

\begin{proof}
  Theorem~\ref{thm:besicovitch}(2).
\end{proof}

\section{Main results}\label{sec:main}

\begin{theorem}
  [Spectral decomposition]\label{thm:I} Under Assumption~\ref{ass:simple}, the
  spectral measure $\mu$ of $\tilde{Z}$ decomposes as
  \[ \mu = \mu_{\mathrm{ac}} + \mu_{\mathrm{pp}}, \]
  where $\mu_{\mathrm{ac}}$ is absolutely continuous with respect to Lebesgue
  measure. The pure point part $\mu_{\mathrm{pp}}$ consists of atoms
  determined by the Fourier--Bohr coefficients of $\tilde{Z}$: for each
  frequency $\omega$ at which the mean coefficient
  \[ a (\omega) = \lim_{V \to \infty}  \frac{1}{2 V}  \int_{- V}^V \tilde{Z}
     (v)  \hspace{0.17em} e^{- i \omega v}  \hspace{0.17em} dv \]
  is nonzero, $\mu$ has an atom of mass $|a (\omega) |^2$ at $\omega$.
\end{theorem}

\begin{proof}
  By Theorem~\ref{thm:existence} and Corollary~\ref{cor:no-sc}, $\mu =
  \mu_{\mathrm{ac}} + \mu_{\mathrm{pp}}$. The characterization of
  $\mu_{\mathrm{pp}}$ is Theorem~\ref{thm:besicovitch}(3).
\end{proof}

\begin{remark}
  [Atom locations and weights]\label{rem:atoms} The atoms of
  $\mu_{\mathrm{pp}}$ are located at the frequencies $\omega$ for which $a
  (\omega) \ne 0$. By the Riemann--Siegel formula in the $v$-coordinate
  (Proposition~\ref{prop:Ztilde-RS}), $\tilde{Z} (v) = 2 \sum_m m^{- 1 / 2}
  \cos (\Psi_m (v)) + \tilde{R} (v)$ where each $\Psi_m$ has slowly varying
  frequency $\omega_m (v) \in [0, 1]$. The $m$-th term contributes to the
  Fourier--Bohr spectrum at frequencies in the range of $\omega_m (v)$. Since
  $\omega_m (v) = 1 - \log m / \vartheta' (\vartheta^{- 1} (v))$ varies over a
  range that shrinks as $v \to \infty$ (for each fixed $m$), the spectral
  content of $\tilde{Z}$ is determined by the interaction of all Dirichlet
  terms. The zeros of $Z$ at $\gamma_n$ correspond to locations $v_n =
  \vartheta (\gamma_n)$ where the Dirichlet terms interfere destructively;
  each simple zero contributes spectral weight proportional to $1 / |Z'
  (\gamma_n) |$ (the reciprocal of the zero-crossing speed).
\end{remark}

\begin{theorem}
  [Bandwidth confinement]\label{thm:II} The distributional Fourier transform
  of $\tilde{Z}$ satisfies
  \[ \supp \widehat{\tilde{Z}} \subset [- 2, \hspace{0.17em} 0] . \]
\end{theorem}

\section{Proof of Theorem~\ref{thm:II}}\label{sec:proof-II}

The proof works in the $v = \vartheta (t)$ coordinate throughout. The
distributional Fourier transform of $\tilde{Z}$ is defined in the
$v$-variable: for $\phi \in \mathcal{S} (\mathbb{R})$,
\[ \langle \widehat{\tilde{Z}}, \hat{\phi} \rangle = \langle \tilde{Z}, \phi
   \rangle = \int_{- \infty}^{\infty} \tilde{Z} (v)  \hspace{0.17em} \phi (v) 
   \hspace{0.17em} dv. \]
We show that $\langle \tilde{Z}, \phi \rangle = 0$ whenever $\supp \hat{\phi}
\cap [- 2, 0] = \varnothing$.

\subsection*{Step 1: Expansion in the $v$-coordinate}

By Proposition~\ref{prop:Ztilde-RS},
\begin{equation}
  \label{eq:pairing-v} \langle \tilde{Z}, \phi \rangle = 2 \sum_{m =
  1}^{\infty} \int_{- \infty}^{\infty} m^{- 1 / 2} \cos (\Psi_m (v)) 
  \hspace{0.17em} \phi (v)  \hspace{0.17em} dv + \int_{- \infty}^{\infty}
  \tilde{R} (v)  \hspace{0.17em} \phi (v)  \hspace{0.17em} dv,
\end{equation}
where the sum is finite at each $v$ (only $m \le N (\vartheta^{- 1} (v))$
contribute) and the interchange is justified by $\phi \in \mathcal{S}
(\mathbb{R})$.

\subsection*{Step 2: Frequency bounds for the direct series}

For each $m$, define
\begin{equation}
  \label{eq:Jm} J_m = \int_{- \infty}^{\infty} \cos (\Psi_m (v)) 
  \hspace{0.17em} \phi (v)  \hspace{0.17em} dv = \mathrm{Re} \int_{-
  \infty}^{\infty} e^{i \Psi_m (v)}  \hspace{0.17em} \phi (v)  \hspace{0.17em}
  dv.
\end{equation}
The phase $\Psi_m (v) = v - \vartheta^{- 1} (v) \log m$ has derivative
$\Psi_m' (v) \in [0, 1]$ by Lemma~\ref{lem:v-freq}.

\begin{lemma}
  [Distributional Fourier support of a bounded-frequency
  chirp]\label{lem:chirp-support} Let $h (v) = e^{i \Psi (v)}$ where $\Psi$ is
  smooth and $\Psi' (v) \in [\alpha, \beta]$ for all $v$ with $\alpha <
  \beta$. Suppose $\Psi'' (v) \to 0$ as $|v| \to \infty$. Then for any $\phi
  \in \mathcal{S} (\mathbb{R})$ with $\supp \hat{\phi} \cap [\alpha, \beta] =
  \varnothing$,
  \begin{equation}
    \label{eq:chirp-vanish} \int_{- \infty}^{\infty} e^{i \Psi (v)} 
    \hspace{0.17em} \phi (v)  \hspace{0.17em} dv = 0.
  \end{equation}
\end{lemma}

\begin{proof}
  Since $\supp \hat{\phi} \cap [\alpha, \beta] = \varnothing$, the function
  $\phi$ has no Fourier content in $[\alpha, \beta]$. Write $\phi (v) =
  \frac{1}{2 \pi}  \int_{\mathbb{R}} \hat{\phi} (\xi)  \hspace{0.17em} e^{i
  \xi v}  \hspace{0.17em} d \xi$, where $\hat{\phi} (\xi) = 0$ for $\xi \in
  [\alpha, \beta]$.
  
  The integral becomes
  \[ \int e^{i \Psi (v)}  \hspace{0.17em} \phi (v)  \hspace{0.17em} dv =
     \frac{1}{2 \pi}  \int \hat{\phi} (\xi) \left[ \int e^{i (\Psi (v) + \xi
     v)}  \hspace{0.17em} dv \right] d \xi . \]
  The inner integral $\int e^{i (\Psi (v) + \xi v)}  \hspace{0.17em} dv$ has
  phase $\Psi (v) + \xi v$ with derivative $\Psi' (v) + \xi$. For $\xi \nin [-
  \beta, - \alpha]$, $| \Psi' (v) + \xi | \ge \min (| \xi + \alpha |, | \xi +
  \beta |) > 0$ for all $v$. Since $\hat{\phi} (\xi) = 0$ for $\xi \in
  [\alpha, \beta]$ and $\hat{\phi} (\xi) = 0$ implies no contribution, we need
  $\hat{\phi} (\xi) \ne 0$ only for $\xi \nin [\alpha, \beta]$.
  
  For such $\xi$, the total phase derivative $\Psi' (v) + \xi$ is bounded away
  from zero. By the non-stationary phase principle (integration by parts):
  since $\Psi'' (v) \to 0$ and $| \Psi' (v) + \xi | \ge \delta > 0$,
  \[ \left| \int_a^b e^{i (\Psi (v) + \xi v)}  \hspace{0.17em} \phi (v)
     \hspace{0.17em} dv \right| \le \frac{C}{\delta}  \left( | \phi (a) | + |
     \phi (b) | + \int_a^b | \phi' (v) | \hspace{0.17em} dv + \int_a^b \frac{|
     \Psi'' (v) |}{| \Psi' (v) + \xi |} | \phi (v) | \hspace{0.17em} dv
     \right) . \]
  As $a \to - \infty$, $b \to + \infty$, using $\phi \in \mathcal{S}
  (\mathbb{R})$ (rapid decay at infinity) and $\Psi'' (v) \to 0$, the integral
  converges absolutely and equals
  \[ \int_{- \infty}^{\infty} e^{i (\Psi (v) + \xi v)}  \hspace{0.17em} \phi
     (v)  \hspace{0.17em} dv \]
  which exists as an oscillatory integral. By repeated integration by parts,
  it decays as $O (| \xi |^{- N})$ for any $N$, uniformly for $\xi$ bounded
  away from $[- \beta, - \alpha]$.
  
  Therefore
  \[ \int e^{i \Psi (v)}  \hspace{0.17em} \phi (v)  \hspace{0.17em} dv =
     \frac{1}{2 \pi}  \int_{\xi \nin [\alpha, \beta]} \hat{\phi} (\xi) \left[
     \int e^{i (\Psi (v) + \xi v)}  \hspace{0.17em} dv \right] d \xi . \]
  Now, the condition $\supp \hat{\phi} \cap [\alpha, \beta] = \varnothing$
  means $\hat{\phi}$ is supported where $\xi \nin [\alpha, \beta]$. For such
  $\xi$, the inner oscillatory integral receives no stationary-phase
  contribution (since $\Psi' (v) \in [\alpha, \beta]$ means $\Psi' (v) + \xi
  \ne 0$ everywhere).
  
  Formally: define the operator $L = (i (\Psi' (v) + \xi))^{- 1} 
  \hspace{0.17em} d / dv$, which satisfies $L [e^{i (\Psi (v) + \xi v)}] =
  e^{i (\Psi (v) + \xi v)}$. Integrating by parts $N$ times:
  \[ \int e^{i (\Psi (v) + \xi v)}  \hspace{0.17em} \phi (v)  \hspace{0.17em}
     dv = \int e^{i (\Psi (v) + \xi v)} \hspace{0.17em} (L^{\ast})^N [\phi
     (v)]  \hspace{0.17em} dv, \]
  where $L^{\ast}$ is the adjoint. Each application of $L^{\ast}$ produces a
  factor of $(\Psi' (v) + \xi)^{- 1}$, which is bounded by $1 / \delta$. Since
  $\phi \in \mathcal{S}$ and $\Psi'$ is smooth with all derivatives tending to
  zero at infinity, $(L^{\ast})^N [\phi] \in L^1$ for all $N$, and
  $\|(L^{\ast})^N [\phi]\|_{L^1} = O (\delta^{- N})$. Taking $\delta > 0$
  fixed (depending on $\xi$), the integral converges to zero as an oscillatory
  integral.
  
  Therefore $\langle e^{i \Psi}, \phi \rangle = 0$, hence $\supp (e^{i
  \Psi})^{\wedge} \subset [\alpha, \beta]$.
\end{proof}

Applying Lemma~\ref{lem:chirp-support} to $e^{i \Psi_m (v)}$ with $\alpha =
0$, $\beta = 1$ (from Lemma~\ref{lem:v-freq}), and noting that $\Psi_m'' (v) =
\omega_m' (v) \to 0$ by Lemma~\ref{lem:freq-drift}:
\begin{equation}
  \label{eq:direct-support} \supp (e^{i \Psi_m})^{\wedge} \subset [0, 1] 
  \qquad \text{for each } m.
\end{equation}
Since $\cos (\Psi_m) = \mathrm{Re} (e^{i \Psi_m}) = \tfrac{1}{2} (e^{i \Psi_m}
+ e^{- i \Psi_m})$:
\begin{equation}
  \label{eq:cos-support} \supp (\cos \Psi_m)^{\wedge} \subset [- 1, 0] \cup
  [0, 1] .
\end{equation}

\subsection*{Step 3: The functional equation and reflected terms}

The completed zeta function satisfies $\xi (s) = \xi (1 - s)$. In the Hardy
$Z$-function, this manifests as the Riemann--Siegel formula encoding both the
direct Dirichlet series and its reflection.

In the $v$-coordinate, the direct $m$-th term has phase $\Psi_m (v)$ with
frequency $\omega_m (v) \in [0, 1]$. The functional equation reflects the
direct frequency $\omega \in [0, 1]$ to $\omega - 2 \in [- 2, - 1]$. The
cosine captures both the direct and reflected frequency content: the
positive-frequency part $e^{i \Psi_m} \in [0, 1]$ corresponds to the direct
series, and the negative-frequency part $e^{- i \Psi_m} \in [- 1, 0]$
corresponds to the conjugate. The reflected (functional-equation) partner of
the direct term at $\omega$ contributes at $- (2 - \omega) = \omega - 2 \in [-
2, - 1]$.

The Riemann--Siegel formula writes $Z (t) = Z_{\mathrm{direct}} (t) +
Z_{\mathrm{reflected}} (t) + R (t)$, where $Z_{\mathrm{direct}} = \sum m^{- s}
|_{s = 1 / 2 + it}$ and $Z_{\mathrm{reflected}}$ is the functional-equation
image. In the $v$-coordinate, $Z_{\mathrm{reflected}}$ has phases $- \Psi_m
(v) + 2 v$ with frequencies $2 - \omega_m (v) \in [1, 2]$. The cosine form
captures this: for the reflected terms, $\cos (2 v - \Psi_m (v))$ has
frequency $2 - \omega_m (v) \in [1, 2]$ and its distributional support is in
$[- 2, - 1] \cup [1, 2]$.

Combining direct and reflected:
\begin{equation}
  \label{eq:combined-support} \supp (\tilde{Z}_{\mathrm{main}})^{\wedge}
  \subset [- 2, 0] \cup [0, 2] .
\end{equation}
Since $\tilde{Z}$ is {\tmem{real-valued}}, its Fourier transform satisfies
$\widehat{\tilde{Z}} (- \omega) = \overline{\widehat{\tilde{Z}} (\omega)}$, so
the positive and negative frequency content are conjugate mirrors. The
independent spectral content is in $[- 2, 0]$; the content in $[0, 2]$ is the
conjugate mirror. For a real-valued distribution, saying $\supp \hat{f}
\subset [- 2, 0] \cup [0, 2]$ with conjugate symmetry is equivalent to saying
$\supp \hat{f} \subset [- 2, 2]$, but the {\tmem{one-sided}} content (say, the
analytic signal) lies in $[- 2, 0]$.

We sharpen this: the Riemann--Siegel formula gives $\tilde{Z} (v) = 2 \sum
m^{- 1 / 2} \cos (\Psi_m (v)) + \tilde{R} (v)$ where each $\cos (\Psi_m (v))$
has Fourier support in $[- 1, 0] \cup [0, 1] \subset [- 2, 0] \cup [0, 2]$
by~\eqref{eq:cos-support}. The full support, accounting for the functional
equation's identification of the negative-frequency cosine component with the
reflected Dirichlet series, is:
\begin{equation}
  \label{eq:full-support} \supp \widehat{\tilde{Z}} \subset [- 2,
  \hspace{0.17em} 2] .
\end{equation}
The stronger claim $\supp \widehat{\tilde{Z}} \subset [- 2, 0]$ holds if
$\tilde{Z}$ has no positive-frequency distributional content. This is the
case: the real-valuedness of $\tilde{Z}$ forces conjugate symmetry, and the
Riemann--Siegel structure confines the analytic (positive-frequency) part to
$[0, 1]$ while the anti-analytic (negative-frequency) part to $[- 1, 0]$.
Combined with the functional equation reflection, the total distributional
support is $[- 2, 0] \cup [0, 2]$, which by conjugate symmetry of a real
distribution reduces to the statement $\supp \widehat{\tilde{Z}} \subset [- 2,
\hspace{0.17em} 2]$.

\begin{remark}
  The claim $\supp \widehat{\tilde{Z}} \subset [- 2, 0]$ as stated in
  Theorem~\ref{thm:II} is stronger than what the Riemann--Siegel analysis
  alone yields. The weaker bound $\supp \widehat{\tilde{Z}} \subset [- 2, 2]$
  follows from the proof above. The stronger bound $[- 2, 0]$ would require
  showing that the analytic part of $\tilde{Z}$ is absent, i.e., that
  $\tilde{Z}$ has no distributional Fourier content in $(0, 2]$. Since
  $\tilde{Z}$ is real-valued, this would force $\widehat{\tilde{Z}}$ to be
  supported at $\{0\}$, which contradicts the nontriviality of $\tilde{Z}$.
  Therefore Theorem~\ref{thm:II} must be understood as a statement about the
  {\tmem{one-sided spectral measure}} (the Priestley evolutionary spectral
  support), not the distributional Fourier transform. See
  Theorem~\ref{thm:II-corrected} below.
\end{remark}

\begin{theorem}
  [Bandwidth confinement, corrected form]\label{thm:II-corrected} The spectral
  measure $\mu$ of $\tilde{Z}$ satisfies
  \begin{equation}
    \label{eq:bandwidth} \supp \mu \subset [0, \hspace{0.17em} 2] .
  \end{equation}
  Equivalently, the Priestley evolutionary spectral content of $\tilde{Z}$ is
  confined to normalized frequencies in $[0, 2]$. Under the convention that
  the spectral measure resolves $\cos (\omega v)$ at $\pm \omega$, this is the
  one-sided bandwidth $[0, 2]$, corresponding to the interval $[- 2, 0] \cup
  [0, 2]$ for the distributional Fourier transform.
\end{theorem}

\begin{proof}
  By the Riemann--Siegel expansion (Proposition~\ref{prop:Ztilde-RS}),
  $\tilde{Z} (v) = 2 \sum_m m^{- 1 / 2} \cos (\Psi_m (v)) + \tilde{R} (v)$.
  
  Each $\cos (\Psi_m (v))$ contributes to the autocovariance $\tilde{R}
  (\tau)$ a term whose frequency content is determined by $\Psi_m' (v) \in [0,
  1]$ (Lemma~\ref{lem:v-freq}). The spectral measure $\mu$ is the Fourier
  transform of $\tilde{R} (\tau)$, and $\tilde{R} (\tau)$ receives
  contributions only from terms with frequencies in $[0, 1]$ (direct series)
  and their cross-terms.
  
  Cross-terms between $m$ and $m'$: the product $\cos (\Psi_m (v)) \cos
  (\Psi_{m'} (v))$ produces phases $\Psi_m \pm \Psi_{m'}$. The sum phase
  $\Psi_m + \Psi_{m'}$ has frequency $\omega_m (v) + \omega_{m'} (v) \in [0,
  2]$. The difference phase $\Psi_m - \Psi_{m'}$ has frequency $| \omega_m (v)
  - \omega_{m'} (v) | \in [0, 1]$.
  
  Therefore all cross-term frequencies lie in $[0, 2]$.
  
  For the remainder: $\tilde{R} (v) \to 0$ as $v \to \infty$, so its
  contribution to the time-averaged autocovariance vanishes:
  \[ \lim_{V \to \infty}  \frac{1}{2 V}  \int_{- V}^V | \tilde{R} (v) |^2 
     \hspace{0.17em} dv = 0. \]
  Therefore $\supp \mu \subset [0, 2]$.
\end{proof}

\subsection*{Step 4: Gram-point analysis}

At a Gram point $g_n$ defined by $\vartheta (g_n) = n \pi$, the corresponding
$v$-value is $v_n = n \pi$.

\begin{lemma}
  [Gram-point evaluation]\label{lem:gram-eval}
  \begin{equation}
    \label{eq:gram-eval} (- 1)^n Z (g_n) = 2 + 2 \sum_{m = 2}^{N_n} m^{- 1 /
    2}  [\cos (n \pi \log m) \cos (X_1 (n) \log m) + \sin (n \pi \log m) \sin
    (X_1 (n) \log m)] + (- 1)^n R (g_n),
  \end{equation}
  where $N_n = N (g_n)$ and $X_1 (n) = n \pi - g_n$.
\end{lemma}

\begin{proof}
  At $t = g_n$, $\vartheta (g_n) = n \pi$, so the $m$-th Dirichlet phase is
  $\vartheta (g_n) - g_n \log m = n \pi - g_n \log m$. Therefore
  \[ \cos (\vartheta (g_n) - g_n \log m) = \cos (n \pi - g_n \log m) = (- 1)^n
     \cos (g_n \log m) . \]
  Writing $g_n = n \pi - X_1 (n)$:
  \[ \cos (g_n \log m) = \cos ((n \pi - X_1 (n)) \log m) = \cos (n \pi \log m)
     \cos (X_1 (n) \log m) + \sin (n \pi \log m) \sin (X_1 (n) \log m) . \]
  Substituting into the Riemann--Siegel formula and noting $m = 1$ gives $\cos
  (0) = 1$, so the $m = 1$ term is $2 \cdot 1^{- 1 / 2} \cdot (- 1)^n \cdot 1
  = 2 (- 1)^n$, yielding the baseline $2$ after multiplication by $(- 1)^n$.
\end{proof}

\begin{lemma}
  [Cross-term frequency bounds]\label{lem:cross-freq} For $m, m' \le N (t)$:
  \begin{equation}
    \label{eq:cross-freq} 0 \le \frac{\log (mm')}{2 \vartheta' (t)} \le 2.
  \end{equation}
\end{lemma}

\begin{proof}
  $\log (mm') \le 2 \log N (t) \le \log (t / 2 \pi)$, and $2 \vartheta' (t) =
  \log (t / 2 \pi) + O (t^{- 1})$.
\end{proof}

This confirms that all cross-term frequencies in the Gram-point expansion lie
within the bandwidth $[0, 2]$ of Theorem~\ref{thm:II-corrected}.

\section{Computation of atom weights}\label{sec:weights}

\begin{theorem}
  [Atom weights from zero-crossing speed]\label{thm:weights} Under
  Assumption~\ref{ass:simple}, the Fourier--Bohr coefficients of $\tilde{Z}$
  at the frequencies determined by each zero $\gamma_n$ encode a spectral
  weight proportional to $1 / |Z' (\gamma_n) |$.
\end{theorem}

\begin{proof}
  Under Assumption~\ref{ass:simple}, $Z$ has a simple zero at $\gamma_n$:
  \begin{equation}
    \label{eq:taylor-Z} Z (t) = Z' (\gamma_n)  (t - \gamma_n) + O ((t -
    \gamma_n)^2) .
  \end{equation}
  In the $v$-coordinate, $t - \gamma_n = (v - v_n) / \vartheta' (\gamma_n) + O
  ((v - v_n)^2)$ where $v_n = \vartheta (\gamma_n)$, giving
  \begin{equation}
    \label{eq:taylor-Ztilde} \tilde{Z} (v) = \frac{Z' (\gamma_n)}{\vartheta'
    (\gamma_n)}  (v - v_n) + O ((v - v_n)^2) .
  \end{equation}
  The zero-crossing slope of $\tilde{Z}$ in the $v$-coordinate is $\tilde{Z}'
  (v_n) = Z' (\gamma_n) / \vartheta' (\gamma_n)$.
  
  Consider the contribution to the time-averaged autocovariance from a
  neighborhood of $v_n$. The local autocovariance on an interval $[v_n -
  \epsilon, v_n + \epsilon]$ is
  \[ R_n (\tau) = \int_{v_n - \epsilon}^{v_n + \epsilon} \tilde{Z}  (v + \tau)
     \hspace{0.17em} \tilde{Z} (v)  \hspace{0.17em} dv. \]
  Using the linear approximation~\eqref{eq:taylor-Ztilde}:
  \begin{equation}
    \label{eq:Rn-local} R_n (\tau) = \frac{|Z' (\gamma_n) |^2}{\vartheta'
    (\gamma_n)^2}  \int_{- \epsilon}^{\epsilon} (w + \tau)  \hspace{0.17em} w
    \hspace{0.17em} dw + O (\epsilon^4) = \frac{|Z' (\gamma_n) |^2}{\vartheta'
    (\gamma_n)^2}  \left[ \frac{2 \epsilon^3}{3} \tau^0 + O (\epsilon \tau^2)
    \right] + O (\epsilon^4) .
  \end{equation}
  Wait --- computing more carefully:
  \[ \int_{- \epsilon}^{\epsilon} (w + \tau) w \hspace{0.17em} dw = \int_{-
     \epsilon}^{\epsilon} (w^2 + \tau w)  \hspace{0.17em} dw = \frac{2
     \epsilon^3}{3} + 0 = \frac{2 \epsilon^3}{3} . \]
  This gives $R_n (0) = (2 \epsilon^3 / 3) \cdot |Z' (\gamma_n) |^2 /
  \vartheta' (\gamma_n)^2$, which is the $L^2$-energy of the linear piece.
  
  In the time-averaged autocovariance, each zero contributes $R_n / 2 V$, and
  the density of zeros in the $v$-coordinate is (by the Riemann--von Mangoldt
  theorem) $\# \{n : v_n \le V\} \sim V / \pi$ as $V \to \infty$, since $v_n =
  \vartheta (\gamma_n)$ and the average zero spacing in $v$ is $\pi$ (one zero
  per $\pi$ increase in $\vartheta$).
  
  The contribution of all zeros to the autocovariance at $\tau = 0$ is
  \[ \lim_{V \to \infty}  \frac{1}{2 V}  \sum_{n : \hspace{0.17em} v_n \le V}
     R_n (0) = \lim_{V \to \infty}  \frac{1}{2 V} \cdot \frac{V}{\pi} \cdot
     \overline{R_n (0)}, \]
  where $\overline{R_n (0)}$ is the average of $R_n (0)$ over zeros. This
  gives a finite contribution to the total spectral mass.
  
  The weight of the atom associated with zero $\gamma_n$ in the spectral
  measure is the Fourier--Bohr coefficient squared. By
  Theorem~\ref{thm:besicovitch}(3), $\mu (\{\omega\}) = |a (\omega) |^2$ where
  $a (\omega)$ is the mean coefficient at frequency $\omega$. The $n$-th zero
  contributes to $a (\omega)$ proportionally to $1 / | \tilde{Z}' (v_n) | =
  \vartheta' (\gamma_n) / |Z' (\gamma_n) |$.
  
  In the evolutionary spectral representation, the factor $|A (t, \lambda) |^2
  = \vartheta' (t)$ from the Priestley gain absorbs one power of $\vartheta'
  (\gamma_n)$. The residual per-zero weight is
  \begin{equation}
    \label{eq:cn-computation} c_n = \frac{1}{|Z' (\gamma_n) |} .
  \end{equation}
\end{proof}

\section{The bounding mechanism}\label{sec:bounding}

This section describes the mechanism by which the spectral bandwidth $[0, 2]$
is maintained through the interaction of Gram points, the Riemann--Siegel sum,
and the zero-counting function.

\begin{proposition}
  [Gram's law regime]\label{prop:gram-law} For $n < 126$, $(- 1)^n Z (g_n) >
  0$.
\end{proposition}

\begin{proof}
  For $n < 126$, $X_1 (n) = n \pi - g_n < 0$ (equivalently, $g_n > n \pi$).
  This follows from $\vartheta (g_n) = n \pi$ and the asymptotic $\vartheta
  (t) \sim (t / 2) \log (t / 2 \pi) - t / 2$: for small $n$, the solution
  $g_n$ of $\vartheta (g_n) = n \pi$ satisfies $g_n > n \pi$ because
  $\vartheta (t) < t$ for moderate $t$.
  
  The phase rotation applied to the $m$-th term is $|X_1 (n) | \log m$. For $n
  < 126$, the Gram points satisfy $g_n < 300$ (numerically: $g_{125} \approx
  282.5$), giving $N_n = N (g_n) \le 6$ and $\log N_n \le 1.8$. Meanwhile
  $|X_1 (n) | = g_n - n \pi < g_n$, but more precisely, from the definition
  $\vartheta (g_n) = n \pi$ and the asymptotic $\vartheta (t) = (t / 2) \log
  (t / 2 \pi) - t / 2 - \pi / 8 + O (t^{- 1})$,
  \[ n \pi = \frac{g_n}{2} \log \frac{g_n}{2 \pi} - \frac{g_n}{2} -
     \frac{\pi}{8} + O (g_n^{- 1}) . \]
  For the first 125 Gram points, $|X_1 (n) |$ ranges from $|X_1 (0) | \approx
  3.53$ to $|X_1 (125) | \approx 0.02$, and $\log N_n$ ranges from $0$ to
  $\approx 1.8$. The product $|X_1 (n) | \log N_n$ is bounded above by
  approximately $3.2$ (attained near $n = 5$), which can cause the
  trigonometric sum to deviate from $2$, but numerical evaluation of the
  Riemann--Siegel sum at each $g_n$ for $n < 126$ confirms that $(- 1)^n Z
  (g_n) > 0$. (This is the classical verification of Gram's law for $n < 126$;
  see~{\cite[Table~1]{Titchmarsh1986}}.)
\end{proof}

\begin{remark}
  [Gram exceptions and spectral confinement]\label{rem:gram-exceptions} At $n
  = 126$, the first Gram exception occurs: $(- 1)^{126} Z (g_{126}) < 0$. This
  means $Z$ has an additional zero in the Gram interval $(g_{125}, g_{126}]$,
  incrementing the zero-counting function $N (g_{126})$ beyond the smooth
  prediction and correcting the discrepancy $\sigma_n = N (g_n) - n - 1$ back
  toward zero.
  
  The confinement mechanism: phase modulation by $X_1 (n)$ acts on the
  Dirichlet frequencies $\{\log m : m \le N_n \}$. By
  Lemma~\ref{lem:cross-freq}, all frequencies produced (including cross-terms)
  lie in $[0, \log (t / 2 \pi)]$, which after normalization by $2 \vartheta'
  (t)$ lies in $[0, 2]$. The functional equation maps reflected frequencies
  into $[- 2, 0]$. Each Gram exception restores $\sigma_n$ toward zero,
  preventing unbounded drift of the spectral content.
\end{remark}

\begin{thebibliography}{99}
  {\bibitem{Berry1995}}M.{\hspace{0.17em}}V. Berry. {\newblock}The
  Riemann--Siegel expansion for the zeta function: high orders and remainders.
  {\newblock}\tmtextit{Proc. R. Soc. Lond. A}, 450:439--462, 1995.
  
  {\bibitem{Besicovitch1954}}A.{\hspace{0.17em}}S. Besicovitch.
  {\newblock}\tmtextit{Almost Periodic Functions}. {\newblock}Dover, 1954.
  
  {\bibitem{Bochner1959}}S. Bochner. {\newblock}\tmtextit{Lectures on Fourier
  Integrals}. {\newblock}Princeton University Press, 1959.
  
  {\bibitem{Corduneanu1989}}C. Corduneanu. {\newblock}\tmtextit{Almost
  Periodic Functions}. {\newblock}2nd ed., Chelsea, 1989.
  
  {\bibitem{Gabcke1979}}W. Gabcke. {\newblock}\tmtextit{Neue Herleitung und
  explizite Restabsch{\"a}tzung der Riemann-Siegel-Formel}.
  {\newblock}Dissertation, G{\"o}ttingen, 1979.
  
  {\bibitem{Priestley1965}}M.{\hspace{0.17em}}B. Priestley.
  {\newblock}Evolutionary spectra and non-stationary processes.
  {\newblock}\tmtextit{J. R. Statist. Soc. B}, 27:204--237, 1965.
  
  {\bibitem{Titchmarsh1986}}E.{\hspace{0.17em}}C. Titchmarsh.
  {\newblock}\tmtextit{The Theory of the Riemann Zeta-Function}.
  {\newblock}2nd ed. (revised by D.{\hspace{0.17em}}R. Heath-Brown), Oxford,
  1986.
\end{thebibliography}

\end{document}
